% Companion Part VI -- VI/33 editorial bridge. Addition only; no existing solution replaced.
\paragraph{Why quasi-coherent sheaves are modules in local coordinates.}
A sheaf of \(\mathcal O_X\)-modules supplies an algebraic module on each open
subset, together with restrictions compatible with the ring of functions.
The key property of a \emph{quasi-coherent} sheaf is that, near any point,
this information comes from an ordinary module over the coordinate ring.
When \(U=\operatorname{Spec}A\), an \(A\)-module \(M\) produces
\(\widetilde M\), characterized on a standard open by
\[
 \Gamma(D(f),\widetilde M)=M_f.
\]
This is a concrete rule rather than a vague local-to-global slogan:
restriction to \(D(f)\) performs the same localization on both the ring
and the module. For example, on \(\mathbb A^1_k=\operatorname{Spec}k[t]\),
\(\widetilde{k[t]/(t)}\) is supported at the origin. Over \(D(t)\) its
sections vanish because \((k[t]/(t))_t=0\).

\paragraph{Kernels, images, and cokernels can be computed on affine opens.}
Take \(A=k[x,y]/(xy)\), let \(M=A\), and consider multiplication by
\(x\):
\[
 \varphi:A\longrightarrow A,\qquad a\longmapsto xa.
\]
Because \(xy=0\), every multiple of \(y\) is killed by \(\varphi\).
Conversely, writing an element of \(A\) as a polynomial in \(x\)
and \(y\) with mixed monomials removed shows that \(xa=0\) implies
\(a\in(y)\). Thus
\[
 \ker\varphi=(y),\qquad \operatorname{im}\varphi=(x),\qquad
 \operatorname{coker}\varphi=A/(x)\cong k[y].
\]
The corresponding morphism of sheaves has kernel
\(\widetilde{(y)}\), image \(\widetilde{(x)}\), and cokernel
\(\widetilde{A/(x)}\). These equalities follow from exactness
of localization. They illustrate precisely why the category of
quasi-coherent sheaves is closed under these operations.

\paragraph{What the Noetherian assumption contributes.}
Quasi-coherence alone does not demand that the associated module be
finitely generated. For a Noetherian ring, submodules of finitely
generated modules are finitely generated. Consequently a homomorphism
between finite modules has a finite kernel and image as well as a finite
cokernel. On a locally Noetherian scheme, such modules are coherent.
The distinction matters: a finite presentation of a map is not
automatically enough to control every kernel over a non-Noetherian base.
In arguments about coherence, explicitly identify the finiteness
hypothesis used at the passage from a module to its submodule.

\paragraph{Pullback changes scalars; pushforward changes the region of observation.}
Suppose \(g:\operatorname{Spec}B\to\operatorname{Spec}A\) is induced
by \(A\to B\). For an \(A\)-module \(M\), the pullback has the
familiar affine description
\[
 g^*(\widetilde M)\cong\widetilde{B\otimes_A M}.
\]
For example, with \(A=k[t]\), \(M=A/(t)\), and
\(B=k[t,u]/(u^2-t)\cong k[u]\), the pulled-back module is
\[
 B\otimes_A A/(t)\cong B/(u^2)=k[u]/(u^2).
\]
The source module was concentrated at \(t=0\); its pullback is a
\emph{nonreduced} fiber of length two. This calculation shows why
pullback does not merely take set-theoretic inverse images.

A pushforward need not retain coherence. For the affine-line structure
map \(p:\mathbb A^1_k\to\operatorname{Spec}k\),
\[
 \Gamma(\operatorname{Spec}k,p_*\mathcal O_{\mathbb A^1_k})
     =k[t],
\]
which is infinite-dimensional as a \(k\)-vector space. Thus
\(p_*\mathcal O_{\mathbb A^1_k}\) is quasi-coherent but not coherent.
Properness and additional finiteness hypotheses are significant in
later coherent pushforward theorems; none should be silently inferred
for this nonproper map.
